\bookchapter{第三十章：投屏、演示與車載互聯}

\chapterintro{手機屏幕再大亦遂7英寸，但汝之電視有55英寸、汝之投影儀有100英寸。將手機畫面“扔”至大屏幕上，體驗乃完全不同維度之。此一章講清楚投屏、演示與車載互聯此三個高頻場景，幫汝將手機之能力延伸至更大之屏幕上。}

\booksection{一、無線投屏三大陣營——勿搞混矣}

無線投屏看起來皆乃“手機畫面至電視上”，但底層技術有三種，互不通用：

\begin{center}
\begin{tabularx}{\linewidth}{|X|X|X|}
\hline
\textbf{技術} & \textbf{誰用} & \textbf{最高速率} \tabularnewline
\hline
AirPlay & iPhone/iPad/Mac & 蘋果生態，體驗最流暢。須Apple TV或支持AirPlay之電視 \tabularnewline
Miracast & 安卓手機 & Wi-Fi Direct直連，不須路由器。安卓手機"無線投屏"默認用此 \tabularnewline
DLNA & 凡平臺 & 唯推送媒體文件（視頻/音樂/圖片），不鏡像屏幕。適合推視頻至電視上播放 \tabularnewline
\hline
\end{tabularx}
\end{center}

\booksubsection{各平臺操法：}

\begin{itemize}
  \item \textbf{iPhone}：控制中心 > 屏幕鏡像 > 選擇汝之電視/Apple TV。若電視不支持AirPlay，可買一個Apple TV盒子或支持AirPlay之電視棒。
  \item \textbf{安卓}：下拉通知欄 > 尋到"無線投屏"/"智能投屏"/"多屏互動"按鈕 > 選擇電視。須手機與電視連同一個Wi-Fi。
  \item \textbf{萬能策}：若電視太老不支持任何無線投屏，買一個Chromecast或小米盒子/華爲盒子插至電視HDMI口上，遂能接收投屏矣。
\end{itemize}

\booksection{二、有線投屏——更穩、零延遲}

無線投屏方便但有延遲，玩遊戲與作演示時體驗不好。有線投屏能解決此問題：

\begin{itemize}
  \item \textbf{iPhone}：須一根Lightning/USB-C轉HDMI之轉換線（約100-200元）。插上後手機畫面直接輸出至電視/投影儀之HDMI口。
  \item \textbf{安卓}：部分手機支持USB-C轉HDMI直連（須手機支持DP Alt Mode，大部分旗艦皆支持）。買一根USB-C to HDMI線即可。
  \item \textbf{會議室場景}：若汝屢須用手機作演示，建議於包裏常備一根轉HDMI線。會議室之電視/投影儀幾皆有HDMI口，插上遂能用，比折騰無線投屏靠譜多矣。
\end{itemize}

\booksection{三、手機變電腦——桌面模式}

部分安卓手機支持“桌面模式”——連上顯示器後，手機會輸出一個類似電腦桌面之界面，可用鼠標鍵盤操作：

\begin{itemize}
  \item \textbf{Samsung DeX}：三星旗艦之招牌功能。連接顯示器後彈出桌面界面，可多窗口操作，體驗最接近真正之電腦。配合藍牙鍵鼠，出差可不帶筆記本。
  \item \textbf{華爲桌面模式}：華爲Mate系列/Pura系列連顯器後亦有類似之能。支多窗口、文管、辦公App。
  \item \textbf{摩託羅拉Ready For}：Moto旗艦亦支持，功能類似DeX。
  \item \textbf{適用場景}：輕辦公（回郵件、改文檔）、演示PPT、大屏觀視頻。不適合重度Excel操作或專業設計。
\end{itemize}

\booksection{四、車載互聯——開車以手機之正法}

開車時操作手機乃違法之，但導航、聽音樂、接電話又乃剛需。車載互聯繫統使汝於車機屏幕上使用手機之核心功能，手機可放於一邊不碰：

\begin{itemize}
  \item \textbf{Apple CarPlay}：iPhone用戶之首選。通過USB或無線連接車機，車機屏幕上出現iPhone之精簡界面——導航（高德/百度）、音樂、電話、信息、播客。Siri語音控制全程免手。
  \item \textbf{Android Auto}：安卓用戶對策。能與CarPlay似，支Google Maps、Spotify等。國內因Google服限，部能不可用。
  \item \textbf{華爲HiCar}：華爲自研方案，國內車型支持率高。除矣導航音樂，還能用華爲視頻通話、泊車影像融合等特色功能。
  \item \textbf{百度CarLife}：跨平臺方案，iOS與安卓皆能用。功能較少但兼容性廣，甚多老車型唯支持CarLife。
\end{itemize}

\begin{tcbwarn}{安全提醒}
無論哪種車載互聯方案，開車時皆應儘量使用語音控制。導航提前設好目的地，音樂提前選好歌單，勿於行駛中操作車機屏幕。

\end{tcbwarn}

\booksection{五、會議室投屏實戰——不復手忙腳亂}

開將乃投屏最高頻之場景，亦乃最容易出問題之場景。以下乃幾條實戰經驗：

\begin{itemize}
  \item \textbf{提前至場測試}：永遠勿等至會議開始方連投屏。提前10分鐘到場，測試好連接方式。Wi-Fi投屏失敗之概率比汝欲之高得多。
  \item \textbf{準備有線方案作Plan B}：於包裏放一根轉HDMI線。無線投屏連不上之時，有線即救命稻草。
  \item \textbf{閉告預覽}：投屏時汝之微信消息將顯示於會議室大屏上。投屏前於設置裏關閉"鎖屏通知內容預覽"，或開啓"演示模式"（部分手機支持，投屏時自動隱藏通知）。
  \item \textbf{用"演示文檔"而非"操作手機"}：若汝要展示PPT，直接以PPT存到手機裏用WPS打開全屏播放，而非投屏整個手機桌面。此樣更專業、更流暢。
\end{itemize}

\begin{tcbpractice}{動手：完成一次手機至電視之投屏}
\begin{enumerate}
  \item 確保手機與電視連同一Wi-Fi。
  \item 以手機開一張照片或一個視頻。
  \item 試無線投屏（iPhone用屏幕鏡像，安卓用無線投屏之鍵）。
  \item 若失敗，檢查電視為否支持對應協議，或嘗試DLNA推送（於相冊/視頻App裏點"投屏"按鈕）。
  \item 投屏成功後，試試於手機上作其他操作，觀大屏為否同步顯示。
\end{enumerate}

\end{tcbpractice}

\begin{tcbsummary}{本章之要}
\begin{itemize}
  \item 無線投屏三陣營：AirPlay（蘋果）、Miracast（安卓）、DLNA（跨平臺推媒體）
  \item 有線投屏更穩零延遲，會議室場景必備轉HDMI線
  \item Samsung DeX/華爲桌面模式使手機變身輕辦公電腦
  \item 車載互聯：CarPlay/Android Auto/HiCar/CarLife四大策
  \item 會議室投屏鐵律：提前測試、準備Plan B、閉告預覽
\end{itemize}

\end{tcbsummary}
